\section{The original comparison cone, its native spectrum, and its arithmetic map}
\label{sec:fci}

This section reconstructs the mathematical intake received on 17 September
2026 from the actual CPQ, SFB and new NTE-threshold coefficient maps.
Here NTE-threshold means \emph{A smaller threshold-profile error from
the original nested flags}.  Its complete source body has SHA256
\begin{center}\small\ttfamily
91f2903b110d83b3e8907b486ccaeee6ef2ccdefa900f4676431b09c3f805b4c
\end{center}
It is retained in \path{providers/NTE_THRESHOLD_COMPLETE.tex}.
It is distinct from the older native-total section which also uses
NTE tags.  The full CPQ and SFB providers are retained beside it.
This section proves the
finite cone and metric identities directly.  The stated remote Wolfram
execution and Lean inspection are not used as evidence.  The complete
period-dependent coefficients remain the original ones, and no spectrum
of an arithmetic Euler operator is identified with a metric Laplacian.

\subsection{Fixed spaces and the original four-copy metric}

Retain the original hypothetical simple quartet, its full unit $v_h=g/h$
for $g=2\xi$, the actual five-orbit period, its branch, and the CPQ
rank-valid domain.  Put $k=4l+1\geq9$, $j=k+4$, $q_a=(a+1)^2$,
$q'_a=(a-7)^2$, and $m_a=16a-48-v$, where
\[
 \begin{gathered}
 E_A(z)=\sum_{a,b=0}^8a_{ab}e^{\beta_{8,a,b}z},\qquad
 v=\operatorname{ord}_0E_A,\quad \mu_d=E_A^{(d)}(0),\quad \mu_v\ne0,\\
 \beta_{a,r,t}=a/2+(2r-a)\delta+\mathrm i(2t-a)\gamma.
 \end{gathered}
 \tag{FCI.1}
\]
All $a_{ab}$ are the original period coefficients.  The two source pairs
remain $(s_0,s_1)=(1,1)$ and $(k,j)$.  At target cutoff
$q_j-1\leq N\leq2q_j$, the positive coefficient metric is
\[
 \mathcal H=\mathbb C^{4m_j},\qquad
 \mathcal D_N=I_4\otimes D_{j,N-v}^{(s_1)}.
 \tag{FCI.2}
\]
The provider notation $\operatorname{diag}(D,4)$ means four copies of
$D$, exactly as its map $\operatorname{diag}(B_j,4)$ has target
$\mathbb C^{4t_j}$.  It does not mean the block matrix $D\oplus[4]$.
The $D$ in FCI.2 includes the full attained minimum over all original
lower evaluations.  Its native affine center is $c'_j=j/2-4$.

Use the original injective $B_{\rm cof}:\mathbb C^{m_k}\to\mathcal H$,
original invariant inclusion $J_0$, original onto
$\mathbb B=I_4\otimes B_j$, and fixed connecting columns $F_e$,
where $e\in\{0,1\}$.  Define the actual spaces
\[
 \begin{gathered}
 I=\operatorname{im}(B_{\rm cof}J_0),\quad
 B=\operatorname{im}B_{\rm cof},\quad
 S=\operatorname{im}[B_{\rm cof},F_e],\quad J=\ker\mathbb B,\\
 K_i=I\cap J,\quad W_i=I+J,\quad W=B+J,\quad
 K_c=S\cap W,\quad T=S+J=S+W.
 \end{gathered}\tag{FCI.3}
\]
Thus $I\subset B\subset K_c\subset S$ and
$J\subset W_i\subset W\subset T$.  These spaces and frames are
fixed as $N$ varies.  The columns of $F_e$ are coordinates on $S/B$;
their number is $r_e=8k+24-3e$.  With the original frame
$C_T=[U_R,U_b,U_o]$, CPQ gives
\[
 \pi_o=[0\ 0\ I]C_T^{-1}\mathbb B,\qquad
 \ker\pi_o=W,\qquad E_e=\pi_oF_e.
 \tag{FCI.4}
\]
Indeed $\pi_oz=0$ is equivalent to $\mathbb Bz$ being in
$\operatorname{im}(\mathbb BB_{\rm cof})$, which is equivalent
to $z-B_{\rm cof}b\in J$ for some $b$.  Onto follows from
the onto map $\mathbb B$ and the full frame $C_T$.
Set $c_e=\operatorname{rank}E_e$ and $h_e=r_e-c_e$.

\subsection{The entire cone, including its two cohomology metrics}

Consider $C=[B\hookrightarrow S]$ and
$D=[W\hookrightarrow\mathcal H]$ in degrees zero and one.
Inclusion is a cochain map, and its induced comparison is
\[
 f:S/B\longrightarrow\mathcal H/W,\qquad [s]_B\longmapsto[s]_W.
 \tag{FCI.5}
\]
Its kernel is $K_c/B$, its image is $T/W$, and its cokernel is
$\mathcal H/T$.  For the image assertion, a value represented by
$s+w\in T$ is the same value as $s$ modulo $W$; the other assertions
follow directly from equality of quotient representatives.
The cone in degrees $-1,0,1$ is the explicit complex
\[
 B\xrightarrow{b\mapsto(b,-b)}W\oplus S
       \xrightarrow{(w,s)\mapsto w+s}\mathcal H.
 \tag{FCI.6}
\]
Its differential squares to zero.  The first differential is injective;
the middle cycles are $(-s,s)$ with $s\in K_c$, modulo pairs from $B$;
the final boundaries are $W+S=T$.  Hence
\[
 H^{-1}=0,\qquad H^0=K_c/B,\qquad H^1=\mathcal H/T.
 \tag{FCI.7}
\]
These identifications are induced by the written maps, not by dimensions.

Use the restricted $\mathcal D_N$ metric in the two end terms and the
orthogonal direct sum of its restrictions in the middle term.  For a
middle cocycle the quotient norm is
\[
 \min_{b\in B}\|(-s,s)+(b,-b)\|^2
       =2\min_{b\in B}\|s-b\|_{\mathcal D_N}^2.
 \tag{FCI.8}
\]
In CPQ coordinates the cohomology metrics are therefore $2G_{h,e}$
on $H^0$ and $Q_{{\rm rem},e}$ on $H^1$.  The factor two is part
of the literal direct-sum metric and is retained before any return.

For clarity, the full native Gram definitions are as follows.  Let
$Z_h$ span $\ker E_e$, let $C_P=[Z_h,Z_c]$ be the fixed invertible
source frame, let $A_e=E_eZ_c$, and choose $R_e^o$ so that
$C_o=[A_e,R_e^o]$ is a fixed invertible target frame.  With $H_e$
the attained Gram on $S/B$ and $Q_o$ the attained Gram on
$\mathcal H/W$ in $\pi_o$ coordinates, put
\[
 \begin{split}
 G_h&=Z_h^*H_eZ_h,\\
 G_s&=Z_c^*H_eZ_c-Z_c^*H_eZ_hG_h^{-1}Z_h^*H_eZ_c,\\
 G_t&=A_e^*Q_oA_e,\\
 Q_{\rm rem}&=(R_e^o)^*Q_oR_e^o
 -(R_e^o)^*Q_oA_eG_t^{-1}A_e^*Q_oR_e^o.
 \end{split}\tag{FCI.9}
\]
Each inverse is on the indicated full-rank value space; an empty
block uses determinant one and its empty inverse.  Completing the
square in each positive block gives the two exact identities
\[
 |\det C_P|^2\det H_e=\det G_h\det G_s,\qquad
 |\det C_o|^2\det Q_o=\det G_t\det Q_{\rm rem}.
 \tag{FCI.10}
\]
The first minimum is over $B$ and the killed-source directions,
hence over $K_c$; the second quotient is over the actual image
$T/W$.  They therefore identify $G_s$ with $S/K_c$, $G_t$ with
$T/W$, and $Q_{\rm rem}$ with $\mathcal H/T$, using the stated
fixed value coordinates.

\subsection{The comparison spectrum and both cone Laplacians}

The actual cohomology comparison has squared singular-value operator
\[
 A_N=H_e^{-1/2}E_e^*Q_oE_eH_e^{-1/2}.
 \tag{FCI.11}
\]
It satisfies $0\preceq A_N\preceq I$: the target norm minimizes
over $W\supset B$, so it is no greater than the source norm on
every connecting vector.  Its kernel has dimension $h_e$.
Restrict the source to the orthogonal complement of this kernel.
The restriction represents the attained source quotient $G_s$;
the same values in the image have Gram $G_t$.  Thus, if its positive
eigenvalues are $\lambda_1,\ldots,\lambda_{c_e}$, then
\[
 0<\lambda_a\leq1,\qquad
 \omega_c=\prod_a\lambda_a=\frac{\det G_t}{\det G_s},\qquad
 \frac{\det(zH_e+E_e^*Q_oE_e)}{\det H_e}
       =z^{h_e}\prod_a(z+\lambda_a).
 \tag{FCI.12}
\]
The analogous target polynomial has zero multiplicity
$\dim(\mathcal H/W)-c_e$.  Its nonzero eigenvalues agree:
write the comparison in orthonormal frames as a rectangular matrix
$M$; $M^*M$ and $MM^*$ have the same positive eigenvalues, because
$M$ and $\lambda^{-1}M^*$ identify the eigenspaces for $\lambda>0$.
Combining FCI.10 gives the complete determinant-line identity
\[
 \omega_c=\frac{|\det C_o|^2}{|\det C_P|^2}
       \frac{\det Q_o\det G_h}{\det H_e\det Q_{\rm rem}}.
 \tag{FCI.13}
\]
Both cohomology metrics and both fixed coordinate changes remain.

For a fixed subspace $X$ with full-rank frame, write
$P_X=X(X^*\mathcal D_NX)^{-1}X^*\mathcal D_N$.
It is the $\mathcal D_N$-orthogonal projector by the normal equations.
The adjoint of the last differential in FCI.6 is
$x\mapsto(P_Wx,P_Sx)$, since taking the inner product against any
$(w,s)$ gives the inner product of $w+s$ against $x$.
Its degree-one Laplacian is therefore
\[
 \Delta_c=P_W+P_S,\qquad \ker\Delta_c=T^{\perp_{\mathcal D_N}}.
 \tag{FCI.14}
\]
The kernel identity follows from
$\langle x,\Delta_cx\rangle=\|P_Wx\|^2+\|P_Sx\|^2$.
The invariant comparison is
$[K_i\hookrightarrow I]\to[J\hookrightarrow W_i]$.
Its cone has the same formula with $(B,S,W,\mathcal H)$ replaced
by $(K_i,I,J,W_i)$.  Its middle and final cohomology both vanish,
because $I\cap J=K_i$ and $I+J=W_i$.  Its degree-one Laplacian,
on precisely $W_i$, is
\[
 \Delta_i=(P_J+P_I)|_{W_i}.
 \tag{FCI.15}
\]
It is positive definite on that space.  Acyclicity alone has made
no assertion about its determinant or about an isometry of quotients.

Here is a complete computation of the positive spectral products.
For two subspaces $X,Y$, let $K=X\cap Y$.  On $K$, $P_X+P_Y=2I$;
on $(X+Y)^\perp$ it is zero.  In the orthogonal complement of $K$,
take a singular-value decomposition of the projection between
$X\ominus K$ and $Y\ominus K$.  A nonzero singular value
$0<\sigma<1$ gives unit vectors $x,y$ with $\langle x,y\rangle=\sigma$.
Put $z=(y-\sigma x)/\sqrt{1-\sigma^2}$.  On the orthonormal
two-plane $(x,z)$ the projector sum is
\[
 \begin{pmatrix}1+\sigma^2&\sigma\sqrt{1-\sigma^2}\\
 \sigma\sqrt{1-\sigma^2}&1-\sigma^2\end{pmatrix},
\]
with eigenvalues $1+\sigma,1-\sigma$.  Values equal to zero
and any unpaired orthogonal directions contribute eigenvalue one.
A singular value one would give an additional vector in $X\cap Y$,
which was removed, so it does not occur here.  The quotient projection
$X/K\to(X+Y)/Y$ has squared singular values $1-\sigma^2$ and
the corresponding unit values.  Multiplying all these eigenvalues proves
\[
 \det{}'\Delta_c=2^{\dim K_c}\omega_c,\qquad
 \det{}'\Delta_i=2^{\dim K_i}\omega_i.
 \tag{FCI.16}
\]
Here the prime multiplies the positive eigenvalues only; the kernel
of FCI.14 still represents the cohomology $\mathcal H/T$.
The invariant ratio is the original $\omega_i=\det H_{\rm inv}/
\det A_{\rm inv}$, with both forms at target order $s_1$.

For the original return
\[
 \mathcal Rf=-f(q_j-1)-f(q_j)+f(2q_j-1)+f(2q_j),\qquad
 \mathcal L_e=\mathcal R\log(\omega_c\omega_i),
 \tag{FCI.17}
\]
all intersection dimensions are fixed.  The constants in FCI.16
therefore cancel because $\mathcal R1=0$, giving exactly
\[
 \mathcal L_e=\mathcal R\left(\log\det{}'\Delta_c+
                                      \log\det{}'\Delta_i\right).
 \tag{FCI.18}
\]
This is a statement about degree-one Laplacians of the finite coefficient
cones; it does not identify them with the arithmetic tensor-sum action.

\subsection{The positive boundary-lift matrices and their finite interval}

Use the fixed connecting complement $Z_c$ of FCI.9 and put
\[
 s_N=(1-P_{K_c})F_eZ_c,\quad t_N=(1-P_W)F_eZ_c,\quad
 b_N=s_N-t_N=(P_W-P_{K_c})F_eZ_c.
 \tag{FCI.19}
\]
Since $K_c\subset W$, the last difference is the orthogonal projector
onto $W\ominus K_c$ applied to these columns.  Thus $b_N$ lies in
that space and $t_N\perp W$.  The normal-equation descriptions
of the two minima give
\[
 G_s=s_N^*\mathcal D_Ns_N,\quad G_t=t_N^*\mathcal D_Nt_N,\quad
 G_s=G_t+b_N^*\mathcal D_Nb_N.
 \tag{FCI.20}
\]
The last identity follows by expanding $s_N=t_N+b_N$; both mixed
terms vanish by the proved orthogonality.  It is an equality of full
matrices, not an estimate on their diagonal entries.
For an original invariant lift frame $L_i$ of $I/K_i$, the identical
calculation with $I,J,K_i$ gives
\[
 b_{i,N}=(P_J-P_{K_i})L_i,\qquad
 A_{\rm inv}=H_{\rm inv}+b_{i,N}^*\mathcal D_Nb_{i,N}.
 \tag{FCI.21}
\]
In particular $A_{\rm inv}$ here is not the independent degree-$k$
proper-source form $Q_{\rm inv}$.

Define the positive boundary-cost operators
\[
 C_c=G_t^{-1/2}b_N^*\mathcal D_Nb_NG_t^{-1/2},\qquad
 C_i=H_{\rm inv}^{-1/2}b_{i,N}^*\mathcal D_Nb_{i,N}H_{\rm inv}^{-1/2}.
 \tag{FCI.22}
\]
Taking determinants after congruence in FCI.20--21 proves
\[
 F_N:=\log\det(I+C_c)+\log\det(I+C_i)
       =-\log(\omega_c\omega_i)\geq0,\qquad
 \mathcal L_e=-\mathcal RF_N.
 \tag{FCI.23}
\]
The vector $b_Nx$ belongs to the comparison fibre $W$; this calculation
does not give it a theta primitive.

If $C\succeq0$ is $c$-square, with trace $u$, its nonnegative
eigenvalues give
\[
 \log(1+u)\leq\log\det(I+C)\leq c\log(1+u/c).
 \tag{FCI.24}
\]
The lower bound follows by expanding the product $\prod(1+\lambda_a)$
and dropping nonnegative terms of degree at least two.  For the upper
bound, concavity of $\log$ applied to the $c$ positive numbers
$1+\lambda_a$ gives their average logarithm at most the logarithm
of their average.  When $c=0$ both bounds mean zero.
The exact traces needed here are
\[
 u_c=\operatorname{Tr}(G_t^{-1}(G_s-G_t)),\qquad
 u_i=\operatorname{Tr}(H_{\rm inv}^{-1}(A_{\rm inv}-H_{\rm inv})).
\]
Let $l_N,u_N$ be the sums of the lower and upper bounds for the
two costs.  The original signs in FCI.23 then give
\[
 \begin{split}
 l_{q_j-1}+l_{q_j}-u_{2q_j-1}-u_{2q_j}
 &\leq\mathcal L_e\\
 &\leq u_{q_j-1}+u_{q_j}-l_{2q_j-1}-l_{2q_j}.
 \end{split}\tag{FCI.25}
\]
This is an explicit finite enclosure in the original data.  Positivity
of the individual costs does not imply a sign of their four-endpoint
difference, and no actual growing coefficient has been evaluated here.

\subsection{The unchanged native rows, threshold profile, and complete receiver}

For each original numerator column, write
$N_a(z)=\sum_{u,w}(Z_{\exp,j})_{uw,a}e^{\beta_{j,u,w}z}$.
The quotient $F_a=N_a/E_A$ is the original analytic germ.  Its exact
derivative recurrence and the exact Gamma row coefficients are
\[
 \begin{split}
 f_{a,m}&=\frac{\nu_{a,m+v}-\sum_{d=v+1}^{m+v}
        \binom{m+v}{d}\mu_df_{a,m+v-d}}
                  {\binom{m+v}{v}\mu_v},\qquad
 \nu_{a,n}=\sum_{u,w}(Z_{\exp,j})_{uw,a}\beta_{j,u,w}^{n},\\
 b_{0,m}&=\mathbf1_{m=0},\quad b_{-1,m}=0,\\
 b_{n+1,m}&=\mathrm i^{-1}(b_{n,m-1}-c'_jb_{n,m})
                       -n(n-1+s_1/2)b_{n-1,m},\\
 h_{n,s_1}&=(2\pi)^{s_1/2}n!(s_1/2)_n,\\
 (t_n)_a&=h_{n,s_1}^{-1/2}\sum_m b_{n,m}f_{a,m},\qquad
 (v_n)_{u,w}=h_{n,s_1}^{-1/2}\sum_m b_{n,m}\beta_{j-8,u,w}^{m}.
 \end{split}\tag{FCI.26}
\]
Here $b_{n,m}=0$ outside $0\leq m\leq n$, so in particular
$b_{n,-1}=0$ in the recurrence.  The first line is the product rule for $E_AF_a=N_a$ solved using
the first nonzero $\mu_v$; the later lines expand the original
three-term Gamma recurrence at the unchanged $c'_j$.
Let $T_0,V_0$ include every preceding numerator and lower row and set
\[
 \begin{split}
 A_L&=V_0^*V_0,\quad
 D_0=T_0^*(I-V_0A_L^{-1}V_0^*)T_0,\\
 \rho_n&=t_n-v_nA_L^{-1}V_0^*T_0,\qquad
 \zeta_n=v_nA_L^{-1}v_n^*,\\
 D_1&=D_0+\frac{\rho_n^*\rho_n}{1+\zeta_n},\qquad
 \mathcal U_n=I_4\otimes\frac{\rho_n}{\sqrt{1+\zeta_n}},\qquad
 \mathcal D_1=\mathcal D_0+\mathcal U_n^*\mathcal U_n.
 \end{split}\tag{FCI.27}
\]
To verify the update, subtract the old attained minimum and write the
change in the lower coefficient as $a$.  The remaining minimum is
$a^*A_La+|\rho_nx-v_na|^2$.  Solving its normal equations, or
completing the square, gives $|\rho_nx|^2/(1+\zeta_n)$.
This proves FCI.27 without deleting a lower-evaluation term.

For every fixed flag $X$ define
\[
 K_{X,n}=\mathcal U_nX(X^*\mathcal D_0X)^{-1}X^*\mathcal U_n^*,\quad
 r_n=\frac{\rho_nD_0^{-1}\rho_n^*}{1+\zeta_n}.
 \tag{FCI.28}
\]
Writing $Y=\mathcal U_n\mathcal D_0^{-1/2}$ gives $YY^*=r_nI_4$
and $K_{X,n}=Y\Pi_XY^*$ for the orthogonal projector onto
$\mathcal D_0^{1/2}X$.  Hence $0\preceq K_{X,n}\preceq r_nI_4$;
nested subspaces give ordered matrices.  The determinant lemma gives
the flag increment $b_{X,n}=\log\det(I_4+K_{X,n})$.
Using the four quotient determinant ratios in FCI.3, the complete
increment of $\log(\omega_c\omega_i)$ is
\[
 b_T-b_W+b_{K_c}-b_S+b_{W_i}-b_J+b_{K_i}-b_I.
 \tag{FCI.29}
\]
This is the original eight-flag expression, with every sign preserved.

For completeness the new NTE-threshold error bound follows directly from these
nested flags.  Put $h(x)=\log(1+x)-\log\max(1,x)=u(x)-v(x)$,
where
\[
 u(x)=\log(1+\min(x,1)),\qquad
 v(x)=\begin{cases}0,&x\leq1,\\\log(2x/(1+x)),&x\geq1.\end{cases}
\]
Both $u,v$ increase from zero.  For an ordered chain of four scalar
values in $[0,r]$, the alternating $h$ sum lies in $[-v(r),u(r)]$:
the disjoint increments of each increasing function lie between zero
and its final value.  The two nested chains
$J\subset W_i\subset W\subset T$ and $K_i\subset I\subset K_c\subset S$
give ordered eigenvalues at each of the four indices, by the minimum
and maximum characterization of ordered Hermitian eigenvalues.
No common eigenbasis is needed.  Summing the scalar bound over those
four indices and subtracting the two chains bounds the row error by
$4(u(r_n)+v(r_n))\leq8\log2$.
The original consecutive weights are one at the two ends and two
inside, with sum $2q_j$.  Thus for the original threshold profile
$\mathcal P_{k,e}$,
\[
 |\mathcal L_e-\mathcal P_{k,e}|
 \leq\sum_n\vartheta_n4(u(r_n)+v(r_n))\leq16q_j\log2,
 \quad |\mathcal P_{k,e}+\mathcal RF_N|\leq16q_j\log2.
 \tag{FCI.30}
\]
At a zero row every matrix and its error are zero.

Write $V_X=\mathcal R\log\det X$ for each complete original
quotient metric.  Fixed-frame Schur elimination gives
$V_o=V_t+V_{\rm rem}$ and
$4S_{K,j}=V_I+V_b+V_o$.  The loss is
$\mathcal L_e=V_I+V_t-V_s-V_A$.  Substitution proves the full receiver
\[
 4S_{K,j}=\mathcal L_e+V_b+V_{\rm rem}+V_s+V_A,
 \qquad S_{R,j}=\mathcal T_{j,s_1}-S_{K,j}.
 \tag{FCI.31}
\]
Every $V_X$ is nonnegative: the native row update increases the form,
and minimization over the same full affine fibre preserves this
order; the return is the positive weighted sum of consecutive
log-determinant increments.  This argument does not set any of these
returns to zero.  An interval for $\mathcal L_e$ must be transported
by the complete affine map FCI.31, with the same correlated values.
At the complete arithmetic receiver, write the retained identity as
\[
 \begin{split}
 F_E^{(k)}&=F_K^{(k)}+F_B^{(k)},\quad
 \Psi_k=F_K^{(k)}+F_E^{\rm ar}-F_E^{(k)},\\
 C_k^{\rm budget}&=F_B^{(k)}+\mathcal W_k^{\rm ar}-\mathcal P_--\mathcal P_+,\qquad
 J_k^{\rm action}=\Psi_k+C_k^{\rm budget}
  =F_E^{\rm ar}+\mathcal W_k^{\rm ar}-\mathcal P_--\mathcal P_+.
 \end{split}\tag{FCI.32}
\]
The cancellation is exact substitution.  Improving an allocation in
FCI.31 cannot be counted once more while freezing its matching
term in FCI.32.

\subsection{An explicit map back to the arithmetic packet, with its actual limits}

There is a concrete coefficient-to-arithmetic map, rather than an
assumed equality of the native cone and the theta complex.  This
subsection stays on CPQ's simple-quartet branch.  Put
$\chi_j(S)=\prod_{a,b=0}^j(S-\beta_{j,a,b})$ and
$E_j=\mathbb C[S]/\chi_j$.  The centers are distinct because
$\delta,\gamma>0$.  For one original coefficient vector $z\in\mathbb C^{m_j}$,
let $c=Z_{\exp,j}z$ and define the original numerator functional
\[
 \ell_z([P])=\sum_{a,b}c_{ab}P(\beta_{j,a,b}).
 \tag{FCI.33}
\]
It is well-defined modulo $\chi_j$ and satisfies
$\ell_z(S^n)=\nu_{z,n}$, the actual numerator derivative in FCI.26.
Before this map the conductor is multiplied back in its original
order: $N_z=E_AF_z$, with
$\nu_{z,n}=\sum_{d=0}^n\binom nd\mu_df_{z,n-d}$.
This is the exact product rule; it retains every $\mu_d$, including
the original leading $\mu_v$ and all phase coefficients.
The native center $j/2-4$ and the numerator center $j/2$ have
therefore not been identified.

Let $p_{ab}$ be the actual degree-less-than-$q_j$ Lagrange remainders
with value one at $\beta_{j,a,b}$ and zero at the other centers.
The residue pairing gives the exact unique polynomial class
\[
 L_jz=\sum_{a,b}\chi_j'(\beta_{j,a,b})c_{ab}p_{ab}\in E_j,
 \qquad
 \ell_z([P])=\sum_\lambda\operatorname{Res}_{S=\lambda}
           \frac{(L_jz)(S)P(S)}{\chi_j(S)}\,dS.
 \tag{FCI.34}
\]
Indeed the residue at a simple root is the value of the numerator
divided by $\chi_j'$; substituting the displayed values proves the
formula term by term.  The pairing is perfect because its evaluation
matrix is diagonal with nonzero entries.  This proves uniqueness and
retains each derivative factor instead of dropping it.
The map $Z_{\exp,j}$ is injective on the original coefficient domain:
if its output is zero, then $N_z=0$, all $f_{z,n}=0$ by FCI.26,
and the native positive form in FCI.2 has value zero on $z$.
Positive definiteness forces $z=0$.  Thus $L_j$ is injective.

Use the original full arithmetic packet source $F_h$, with
$\mathcal MF_h=v_h$ and $h(D)F_h=\Theta\phi_*$.  For $P\in E_j$
define
\[
 \iota_j(P)=[r_{\chi_j}(P)(D_1+\cdots+D_j)F_h^{\otimes j}]
       \in Q^{\otimes j},\qquad Q=\mathscr B/\Theta V.
 \tag{FCI.35}
\]
The original full Taylor map sends this class to
$\upsilon_h^{\otimes j}P(A_h^{(j)})1$.
All ordered root tuples above each sum are retained.  The cyclic
annihilator is exactly $\chi_j$, so this map is injective; multiplication
by the full $\upsilon_h=j_h(g/h)$ is invertible.  Polynomial multiples
of $\chi_j(S_\Sigma)$ have the actual ordered theta primitives proved
by division into the $h(s_i)$, using the original $(-1)^{i-1}$ signs
and the admitted source $W_*$.  Thus FCI.35 is independent of the
chosen representative by a literal theta boundary, as in FDC.15a--b.

Taking four independent copies gives the proved injective map
\[
 \mathfrak A:\mathcal H\longrightarrow(Q^{\otimes j})^{\oplus4},
 \qquad \mathfrak A=\iota_j^{\oplus4}(I_4\otimes L_j).
 \tag{FCI.36}
\]
It carries the entire coefficient cone FCI.6 isomorphically to the
relative cone with terms $\mathfrak A(B)$,
$\mathfrak A(W)\oplus\mathfrak A(S)$, and $\mathfrak A(\mathcal H)$,
using exactly $(b,-b)$ and $w+s$.  All kernels, images and quotients
commute with this injection, since equality of images is equivalent
to equality of coefficient vectors.  In particular it gives the
actual quotient isomorphism
\[
 \mathcal H/W\xrightarrow{\sim}
       \mathfrak A(\mathcal H)/\mathfrak A(W).
 \tag{FCI.37}
\]
This is the exact bridge to a relative quotient inside arithmetic
cohomology.  It is not a map identifying $W$ with original theta
boundaries.  Indeed the $81$ centers $\beta_{8,a,b}$ are distinct.
If $\mu_0,\ldots,\mu_{80}$ all vanished, the Vandermonde matrix
on these centers would force every coefficient $a_{ab}$ to vanish,
contradicting $\mu_v\ne0$.  Consequently $v\leq80$, and for $k\geq9$
one has $m_k=16k-48-v\geq16$.  Injectivity of $B_{\rm cof}$ then gives
$0\ne B\subset W$.  Injectivity in FCI.36 gives
$\mathfrak A(W)\ne0$ in actual arithmetic cohomology.
In particular this specific map $\mathfrak A$ cannot factor through
$\mathcal H/W$ with target $(Q^{\otimes j})^{\oplus4}$, since
such a factorization would kill $W$.  FCI.37 is the strongest
quotient map furnished by this constructed injection; it retains
the relative classes instead of declaring them theta boundaries.

Its original arithmetic metric is also explicit.  Let $M_j$ be the
increasing-power coefficient matrix of $I_4\otimes r_{\chi_j}L_j$.
Let $H_{q_j-1}^{\rm ar}$ be the original monomial moment Gram for
$S=j/2+\mathrm iy$ and density
\[
 m_j^{\rm ar}(y)=w_h^{*j}(y),\qquad
 w_h(y)=|v_h(1/2+\mathrm iy)|^2/(2\pi).
\]
The literal representatives in FCI.35 have pullback form
\[
 \Gamma_{\rm ar}=M_j^*(I_4\otimes H_{q_j-1}^{\rm ar})M_j\succ0.
 \tag{FCI.38}
\]
This follows by expanding their full original source norm; its mass
is $(\int w_h)^j$ in each copy.  The exact relative positive operator
against CPQ's native form is
$\mathcal D_N^{-1/2}\Gamma_{\rm ar}\mathcal D_N^{-1/2}$.
No identity or uniform bound for this operator is asserted.
The formula supplies the actual finite comparison and its full
matrix entries while preserving both centers and the conductor.

Finally let $\mathsf A$ be multiplication by $S$ on $E_j^{\oplus4}$,
and put $M=I_4\otimes L_j$.  With the original positive
arithmetic quotient Gram $G=I_4\otimes G_{q_j-1}^{\rm ar}$ on this
full packet, the exact action
compression and its remainder are
\[
 A_{\mathcal H}=(M^*GM)^{-1}M^*G\mathsf A M,\qquad
 R=\mathsf A M-MA_{\mathcal H},\qquad M^*GR=0.
 \tag{FCI.39}
\]
Every inverse exists by injectivity and positivity.  The last identity
is direct substitution.  The arithmetic map intertwines this action
on $\mathcal H$ exactly when $R=0$; in general its full cohomology
defect is the specified injective packet image $\iota_j^{\oplus4}R$.
To induce it on FCI.6 one additionally needs preservation of its actual
subspaces $B,S,W$.  Formula FCI.39 provides no such unproved preservation.
Its remainder is an explicit matrix on the original data, not a
discarded scalar error.  Thus the native Laplacian identities above
do not acquire an arithmetic action or Deligne weights by terminology.

Every construction acts on each original support fibre by retaining its
outer source label, leg mask, and coefficient face; it maps a killed
vector to the supported zero of that target fibre and maps external
$\tau$ to external $\tau$.  The independent-variable connecting maps
of FDC.13--15 remain available on the theta side.  The present finite
cone is related to them by the actual packet maps FCI.33--37, with
its relative arithmetic subspaces explicitly retained.  None of the
calculated finite identities excludes an actual offcritical packet or
certifies such a zero.
